% Generated draft: requires mathematical and source review.
\documentclass[11pt]{article}
\usepackage[margin=25mm]{geometry}
\usepackage{fontspec}
\setmainfont{Latin Modern Roman}
\usepackage{amsmath,amssymb,mathtools,unicode-math}
\setmathfont{Latin Modern Math}
\usepackage{enumitem,needspace,xurl,hyperref}
\setlength{\parindent}{0pt}
\setlength{\parskip}{5pt}
\setlength{\emergencystretch}{3em}
\newcommand{\N}{\mathbb{N}}
\newcommand{\Z}{\mathbb{Z}}
\newcommand{\Q}{\mathbb{Q}}
\newcommand{\R}{\mathbb{R}}
\newcommand{\C}{\mathbb{C}}
\newcommand{\F}{\mathbb{F}}
\newcommand{\A}{\mathbb{A}}
\newcommand{\PP}{\mathbb{P}}
\newcommand{\E}{\mathbb{E}}
\newcommand{\eps}{\varepsilon}
\newcommand{\dd}{\,\mathrm{d}}
\begin{document}
\section*{UnsolvedMath ID: 30003638}
\textbf{Source problem code:} catalogue entry
\subsection{Definitions and mathematical statement}
Under the usual discrete Beurling convention, a generalized prime system is a nondecreasing sequence \(\mathcal P=(p_j)_{j\geq 1}\) of real numbers with \(p_j>1\) and \(p_j\to\infty\). Regard its generalized integers as formal products \(n=\prod_j p_j^{\nu_j}\), where \(\nu_j\in\Z_{\geq0}\) and only finitely many exponents are nonzero. Write \(|n|=\prod_j p_j^{\nu_j}\) for the numerical value. Distinct formal products are counted separately even when their numerical values coincide. Define
\[
\mu_{\mathcal P}(n)=
\begin{cases}
(-1)^{\sum_j\nu_j},&\nu_j\in\{0,1\}\text{ for every }j,\\
0,&\text{otherwise},
\end{cases}
\]
and, for real \(x\geq1\),
\[
N_{\mathcal P}(x)=\sum_{|n|\leq x}1,
\qquad
M_{\mathcal P}(x)=\sum_{|n|\leq x}\mu_{\mathcal P}(n).
\]
The empty product has value \(1\) and Möbius value \(1\).

\textbf{First question, under this convention.} Does there exist a generalized prime system \(\mathcal P\) and a real number \(\alpha<\tfrac12\) such that
\[
\exists C>0\ \exists x_0\geq1\ \forall x\in\R,\quad
x\geq x_0\ \Longrightarrow\ |M_{\mathcal P}(x)|\leq Cx^\alpha?
\]
The supplied record does not specify a density hypothesis, an asymptotic normalization for \(N_{\mathcal P}\), or further admissibility conditions. In particular, the reference to periodic errors \(N(x)-x\) does not state whether a condition such as \(N_{\mathcal P}(x)\sim x\) is required in this existence question. The definitions above make the standard discrete interpretation explicit; they do not establish that it captures every hypothesis of the intended report question.

\textbf{Second question, requiring clarification.} Can the construction method involving periodic errors \(N(x)-x\) be adapted to better nonperiodic approximations of \(\lfloor x\rfloor\)? The original statement further requests approximations whose zero real parts have supremum greater than \(\tfrac12\). If \(Z\) denotes the intended zero set, that wording asks for
\[
\sup\{\operatorname{Re}\rho:\rho\in Z\}>\tfrac12.
\]
The cleaned statement instead describes associated Mellin transforms with zeros whose real parts approach \(\tfrac12\). It specifies neither the side of approach nor the supremum of those real parts. These descriptions do not specify an equivalent zero condition. Moreover, the record gives neither the approximation class and meaning of ``better'', nor the construction method, Mellin-transform formula, analytic continuation, or intended zero set. Consequently this second question cannot be made into a unique quantified mathematical statement from the supplied evidence.
\subsection{Short English statement}
The first question asks whether a Beurling generalized prime system can have summatory Möbius function \(O(x^\alpha)\) for some real \(\alpha<1/2\). Its intended additional hypotheses are unspecified. A related question concerns adapting periodic constructions to nonperiodic approximations of \(\lfloor x\rfloor\), but the approximation criterion, transform, and required zero condition need clarification. The supplied open label is a dated curation claim.
\subsection{Sources}
\begin{itemize}
\item Supplied local registry source record: permanent UnsolvedMath ID \(30003638\), exact source problem code \texttt{catalogue entry}, title ``Sub-Square-Root Möbius Bounds for Generalized Primes'', in the Oberwolfach Reports Open Problems collection. No individual proposer is supplied; the proposed year is recorded as 2017.
\item Supplied source citation: \emph{Mini-Workshop: Interplay between Number Theory and Analysis for Dirichlet Series} (2018), Oberwolfach Reports, pp.\ 3035--3069, DOI 10.4171/owr/2017/51. \url{https://doi.org/10.4171/owr/2017/51}. The background describes these reports as records of workshop discussions and open problems.
\item Supplied literature source: ``Oberwolfach Report 51, open-problems session''. \url{https://doi.org/10.4171/owr/2017/51}.
\item Supplied literature source: Hilberdink, \emph{Well-behaved Beurling primes and integers}. \url{https://doi.org/10.1016/j.jnt.2004.09.001}.
\item The literature assessment dated exactly \texttt{2026-08-21} is supplied dated curation, not independently verified here. It labels the problem open and states that the nonperiodic question of improving the exponent below \(1/2\) remains unresolved, while later work gives related constructions and \(\Omega\)-results but no general affirmative theorem. The supplied background expressly requests independent verification of time-sensitive status and cited claims.
\item The registry labels the statement \texttt{corrected\_verified} and reports checking it against the official report. This is a supplied verification claim, not an independent verification in this draft; it does not resolve the differences between the original and cleaned formulations.
\end{itemize}
\end{document}
